\documentclass[10pt,a4paper]{amsart}
\usepackage{amsmath,amssymb,amsthm,graphicx}
% Retained minimal preamble for the Stacks Project chapter more-morphisms.tex (More on Morphisms).
% Verbatim excerpt of the chapter preamble: Xy-pic setup (native preamble lines 14-18) and the native
% theorem declarations with the mathematical macros (native preamble lines 154-217). Copied mechanically;
% nothing added, removed or redefined. The wrapper class and the amsmath/amssymb/amsthm/graphicx base
% packages are supplied by the assembler, so the class/externaldocument part of the native preamble is omitted.
% For commutative diagrams we use Xy-pic
\usepackage[all]{xy}

% We use 2cell for 2-commutative diagrams.
\xyoption{2cell}
% Theorem environments.
%
\theoremstyle{plain}
\newtheorem{theorem}[subsection]{Theorem}
\newtheorem{proposition}[subsection]{Proposition}
\newtheorem{lemma}[subsection]{Lemma}

\theoremstyle{definition}
\newtheorem{definition}[subsection]{Definition}
\newtheorem{example}[subsection]{Example}
\newtheorem{exercise}[subsection]{Exercise}
\newtheorem{situation}[subsection]{Situation}

\theoremstyle{remark}
\newtheorem{remark}[subsection]{Remark}
\newtheorem{remarks}[subsection]{Remarks}

\numberwithin{equation}{subsection}

% Macros
%
\def\lim{\mathop{\mathrm{lim}}\nolimits}
\def\colim{\mathop{\mathrm{colim}}\nolimits}
\def\Spec{\mathop{\mathrm{Spec}}}
\def\Hom{\mathop{\mathrm{Hom}}\nolimits}
\def\Ext{\mathop{\mathrm{Ext}}\nolimits}
\def\SheafHom{\mathop{\mathcal{H}\!\mathit{om}}\nolimits}
\def\SheafExt{\mathop{\mathcal{E}\!\mathit{xt}}\nolimits}
\def\Sch{\mathit{Sch}}
\def\Mor{\mathop{\mathrm{Mor}}\nolimits}
\def\Ob{\mathop{\mathrm{Ob}}\nolimits}
\def\Sh{\mathop{\mathit{Sh}}\nolimits}
\def\NL{\mathop{N\!L}\nolimits}
\def\CH{\mathop{\mathrm{CH}}\nolimits}
\def\proetale{{pro\text{-}\acute{e}tale}}
\def\etale{{\acute{e}tale}}
\def\QCoh{\mathit{QCoh}}
\def\Ker{\mathop{\mathrm{Ker}}}
\def\Im{\mathop{\mathrm{Im}}}
\def\Coker{\mathop{\mathrm{Coker}}}
\def\Coim{\mathop{\mathrm{Coim}}}

% Boxtimes
%
\DeclareMathSymbol{\boxtimes}{\mathbin}{AMSa}{"02}

%
% Macros for moduli stacks/spaces
%
\def\QCohstack{\mathcal{QC}\!\mathit{oh}}
\def\Cohstack{\mathcal{C}\!\mathit{oh}}
\def\Spacesstack{\mathcal{S}\!\mathit{paces}}
\def\Quotfunctor{\mathrm{Quot}}
\def\Hilbfunctor{\mathrm{Hilb}}
\def\Curvesstack{\mathcal{C}\!\mathit{urves}}
\def\Polarizedstack{\mathcal{P}\!\mathit{olarized}}
\def\Complexesstack{\mathcal{C}\!\mathit{omplexes}}
% \Pic is the operator that assigns to X its picard group, usage \Pic(X)
% \Picardstack_{X/B} denotes the Picard stack of X over B
% \Picardfunctor_{X/B} denotes the Picard functor of X over B
\def\Pic{\mathop{\mathrm{Pic}}\nolimits}
\def\Picardstack{\mathcal{P}\!\mathit{ic}}
\def\Picardfunctor{\mathrm{Pic}}
\def\Deformationcategory{\mathcal{D}\!\mathit{ef}}
\usepackage{cleveref}
\crefname{originalreference}{original reference}{original references}
\Crefname{originalreference}{Original reference}{Original references}
\makeatletter
\newlabel{algebra-lemma-pullback-module}{{02M1}{source}{}{sourcepointer-1}{}}
\newlabel{algebra-lemma-pullback-module@cref}{{[originalreference][0][]02M1}{[1][0][]source}{}{}{}}
\makeatother
\crefname{originalreference}{original reference}{original references}
\Crefname{originalreference}{Original reference}{Original references}
\makeatletter
\newlabel{algebra-lemma-pushdown-module}{{02M0}{source}{}{sourcepointer-2}{}}
\newlabel{algebra-lemma-pushdown-module@cref}{{[originalreference][0][]02M0}{[1][0][]source}{}{}{}}
\makeatother
\crefname{originaldefinition}{original definition}{original definitions}
\Crefname{originaldefinition}{Original definition}{Original definitions}
\makeatletter
\newlabel{definition-weighting}{{0F3A}{source}{}{sourcepointer-3}{}}
\newlabel{definition-weighting@cref}{{[originaldefinition][0][]0F3A}{[1][0][]source}{}{}{}}
\makeatother
\crefname{originalreference}{original reference}{original references}
\Crefname{originalreference}{Original reference}{Original references}
\makeatletter
\newlabel{fields-definition-insep-degree}{{030L}{source}{}{sourcepointer-4}{}}
\newlabel{fields-definition-insep-degree@cref}{{[originalreference][0][]030L}{[1][0][]source}{}{}{}}
\makeatother
\crefname{originalreference}{original reference}{original references}
\Crefname{originalreference}{Original reference}{Original references}
\makeatletter
\newlabel{morphisms-lemma-finite-flat}{{02KB}{source}{}{sourcepointer-5}{}}
\newlabel{morphisms-lemma-finite-flat@cref}{{[originalreference][0][]02KB}{[1][0][]source}{}{}{}}
\makeatother
\title{Length times inseparable degree is a positive weighting for flat locally quasi-finite morphisms}
\author{The Stacks Project Authors}
\date{stacks-project:more-morphisms.tex:local-snapshot-20261011; GNU FDL 1.2 or later}
\begin{document}
\maketitle
\noindent\textit{Editorial note.} Counted span is the author's own proof environment, source lines 22352-22407, immediately following the statement at 22336-22350.\par\medskip
\noindent\textit{Editorial note.} The displayed align environment computing the integral equality is retained verbatim; it is the heart of the proof.\par\medskip
\noindent\textit{Editorial note.} Four cross-chapter pointers (Algebra, Fields, Morphisms) plus the chapter's own definition are recorded in assembly.external\_references.\par\medskip
\noindent\textit{Allowed standard prerequisites (editorial).}\begin{itemize}
\item Definition of a weighting. chapter definition definition-weighting, cited at proof line 22352 onward. Fixes what has to be verified: the integral of the weighting along every finite piece is locally constant.
\item Inseparable and separable degrees of a field extension. Fields, definition fields-definition-insep-degree, cited in the statement. Supplies the inseparable degree factor in the definition of w.
\item Pushing a module down along a finite flat morphism. Algebra, lemmas algebra-lemma-pullback-module and algebra-lemma-pushdown-module, cited at proof lines 22394-22405. Computes the length of the pushed module in terms of the local lengths and the degree, which is exactly the displayed chain of equalities.
\item Finiteness after an elementary etale neighbourhood. Morphisms, lemma morphisms-lemma-finite-flat, cited in the proof. Makes the base change of f finite and flat over a suitable neighbourhood, so that the pushforward computation applies.
\end{itemize}\medskip
\ifdefined\hypertarget\hypertarget{sourcepointer-1}{}\fi
\noindent\textit{Source reference.} Original reference 02M1: focused-seeds/geometry/S9-morph-weighting-flat-quasi-finite/ref-tags.txt, source lines 4--4. Navigation pointer only: the counted proof cites this statement by \textbackslash{}ref as a standing input; the deduction the counted proof has to make is carried out inside the counted span, and no part of that deduction is delegated to this pointer.\par\medskip
\ifdefined\hypertarget\hypertarget{sourcepointer-2}{}\fi
\noindent\textit{Source reference.} Original reference 02M0: focused-seeds/geometry/S9-morph-weighting-flat-quasi-finite/ref-tags.txt, source lines 5--5. Navigation pointer only: the counted proof cites this statement by \textbackslash{}ref as a standing input; the deduction the counted proof has to make is carried out inside the counted span, and no part of that deduction is delegated to this pointer.\par\medskip
\ifdefined\hypertarget\hypertarget{sourcepointer-3}{}\fi
\noindent\textit{Source reference.} Original definition 0F3A: sources/stacks-more-morphisms/pinned-original-source.tex, source lines 22152--22152. Navigation pointer only: the counted proof cites this statement by \textbackslash{}ref as a standing input; the deduction the counted proof has to make is carried out inside the counted span, and no part of that deduction is delegated to this pointer.\par\medskip
\ifdefined\hypertarget\hypertarget{sourcepointer-4}{}\fi
\noindent\textit{Source reference.} Original reference 030L: focused-seeds/geometry/S9-morph-weighting-flat-quasi-finite/ref-tags.txt, source lines 6--6. Navigation pointer only: the counted proof cites this statement by \textbackslash{}ref as a standing input; the deduction the counted proof has to make is carried out inside the counted span, and no part of that deduction is delegated to this pointer.\par\medskip
\ifdefined\hypertarget\hypertarget{sourcepointer-5}{}\fi
\noindent\textit{Source reference.} Original reference 02KB: focused-seeds/geometry/S9-morph-weighting-flat-quasi-finite/ref-tags.txt, source lines 7--7. Navigation pointer only: the counted proof cites this statement by \textbackslash{}ref as a standing input; the deduction the counted proof has to make is carried out inside the counted span, and no part of that deduction is delegated to this pointer.\par\medskip
\begin{lemma}
\label{lemma-weighting-flat-quasi-finite}
Let $f : X \to Y$ be a morphism of schemes. Assume $f$ is
locally quasi-finite, locally of finite presentation, and flat.
Then there is a positive weighting $w : X \to \mathbf{Z}_{> 0}$ of $f$
given by the rule that sends $x \in X$ lying over $y \in Y$ to
$$
w(x) =
\text{length}_{\mathcal{O}_{X, x}}
(\mathcal{O}_{X, x}/\mathfrak m_y \mathcal{O}_{X, x})
[\kappa(x) : \kappa(y)]_i
$$
where $[\kappa' : \kappa]_i$ is the inseparable degree
(Fields, Definition \ref{fields-definition-insep-degree}).
\end{lemma}

\begin{proof}
Consider a diagram as in Definition \ref{definition-weighting}.
Let $u \in U$ with images $x, y, v$ in $X, Y, V$. Then we claim that
$$
\text{length}_{\mathcal{O}_{X, x}}
(\mathcal{O}_{X, x}/\mathfrak m_y \mathcal{O}_{X, x}) =
\text{length}_{\mathcal{O}_{U, u}}
(\mathcal{O}_{U, u}/\mathfrak m_v \mathcal{O}_{U, u})
$$
and
$$
[\kappa(x) : \kappa(y)]_i =
[\kappa(u) : \kappa(v)]_i
$$
The first equality follows as $\mathcal{O}_{X, x} \to \mathcal{O}_{U, u}$
is a flat local homomorphism such that
$\mathfrak m_y \mathcal{O}_{U, u} = \mathfrak m_v \mathcal{O}_{U, u}$
and $\mathfrak m_x \mathcal{O}_{U, u} = \mathfrak m_u$
(because $\mathcal{O}_{Y, y} \to \mathcal{O}_{V, v}$ and
$\mathcal{O}_{X, x} \to \mathcal{O}_{U, u}$ are unramified)
and hence the equality by Algebra, Lemma \ref{algebra-lemma-pullback-module}.
The second equality follows because $\kappa(v)/\kappa(y)$ is a finite
separable extension and $\kappa(u)$ is a factor of
$\kappa(x) \otimes_{\kappa(y)} \kappa(v)$ and hence the inseparable
degree is unchanged. Having said this, we see that formation
of the function in the lemma commutes with \'etale base change.
This reduces the problem to the discussion of the next paragraph.

\medskip\noindent
Assume that $f$ is a finite, flat morphism of finite presentation.
We have to show that $\int_f w$ is locally constant on $Y$.
In fact, $f$ is finite locally free
(Morphisms, Lemma \ref{morphisms-lemma-finite-flat})
and we will show that $\int_f w$ is equal to the degree of $f$
(which is a locally constant function on $Y$). Namely,
for $y \in Y$ we see that
\begin{align*}
(\textstyle{\int}_f w)(y)
& =
\sum\nolimits_{f(x) = y}
\text{length}_{\mathcal{O}_{X, x}}
(\mathcal{O}_{X, x}/\mathfrak m_y \mathcal{O}_{X, x})
[\kappa(x) : \kappa(y)]_i
[\kappa(x) : \kappa(y)]_s \\
& =
\sum\nolimits_{f(x) = y}
\text{length}_{\mathcal{O}_{X, x}}
(\mathcal{O}_{X, x}/\mathfrak m_y \mathcal{O}_{X, x})
[\kappa(x) : \kappa(y)] \\
& =
\text{length}_{\mathcal{O}_{Y, y}}((f_*\mathcal{O}_X)_y/
\mathfrak m_y (f_*\mathcal{O}_X)_y)
\end{align*}
Last equality by Algebra, Lemma \ref{algebra-lemma-pushdown-module}.
The final number is the rank of $f_*\mathcal{O}_X$ at $y$ as desired.
\end{proof}

\end{document}
